% Uppsala University Thesis Template
% Adapted for Master's Thesis in Computer Science
% Original template for Doctoral Theses
% Ver 5.2 - 2012-08-08                  
%                                         
%%%%%%%%%%%%%%%%%%%%%%%%%%%%%%%%%%%%%%%%%%%


\documentclass{UUThesisTemplate}

\usepackage{fix-cm}

% Package to determine whether XeTeX is used
\usepackage{ifxetex}

\ifxetex
	% XeTeX specific packages and settings
	% Language, diacritics and hyphenation
	\usepackage[english]{babel}

	% Font settings with fontspec for proper shape support
	\usepackage{fontspec}
	\setmainfont{TeX Gyre Termes}
	\setsansfont{TeX Gyre Heros}
	\setmonofont[HyphenChar=45]{TeX Gyre Cursor}
\else
	% Plain LaTeX specific packages and settings
	% Language, diacritics and hyphenation
	\usepackage[swedish,english]{babel}

	% Font settings (UTF-8 is now standard, not latin1)
	\usepackage{type1cm}
	\usepackage[utf8]{inputenc}
	\usepackage[T1]{fontenc}
	\usepackage{mathptmx}
	
	% Enable scaling of images on import
	\usepackage{graphicx}
\fi

% Graphics path for figures
\graphicspath{{figures/}{images/}}

% Allow hyphenation and line breaks in monospace (\texttt) text
\ifxetex
	% For XeTeX: redefine \texttt to allow hyphenation
	\renewcommand{\texttt}[1]{{\ttfamily\hyphenchar\font=45\relax #1}}
\else
	\usepackage[htt]{hyphenat}
\fi

% Allow TeX to stretch whitespace to avoid overfull hboxes
\setlength{\emergencystretch}{3em}

% Code listings for Java/C++ code (essential for CS thesis)
\usepackage{listings}
\usepackage{xcolor}

% Configure listings for Java code
\lstset{
	basicstyle=\ttfamily\small,
	keywordstyle=\color{blue}\bfseries,
	commentstyle=\color{gray}\itshape,
	stringstyle=\color{red},
	numbers=left,
	numberstyle=\tiny\color{gray},
	stepnumber=1,
	numbersep=8pt,
	backgroundcolor=\color{white},
	showspaces=false,
	showstringspaces=false,
	showtabs=false,
	frame=single,
	tabsize=4,
	captionpos=b,
	breaklines=true,
	breakatwhitespace=false,
	escapeinside={(*@}{@*)},
	language=Java
}

% URL handling
\usepackage{url}
\urlstyle{same}


% Tables
\usepackage{booktabs}
\usepackage{tabularx}

% Document links and bookmarks (modern configuration)
\usepackage{hyperref}
\hypersetup{
	colorlinks=true,
	linkcolor=black,
	citecolor=blue,
	urlcolor=blue,
	pdftitle={Optimizing Weak Reference Processing in the JVM Z Garbage Collector},
	pdfauthor={Fredrik Hammarberg},
	pdfsubject={Master's Thesis in Computer Science},
	pdfkeywords={Java, JVM, Garbage Collection, ZGC, Weak References, Performance Optimization}
} 

% Improve PDF bookmarks and avoid rerunfilecheck warnings
% (Load after hyperref and before cleveref)
\usepackage{bookmark}

% Better cross-referencing (must be loaded AFTER hyperref)
\usepackage{cleveref}

% Numbering of headings down to the subsection level
\numberingdepth{subsection}

% Including headings down to the subsection level in contents
\contentsdepth{subsection}


% Define thesis metadata
\title{Optimizing Weak Reference Processing in the JVM Z Garbage Collector}
\subtitle{A Performance Analysis and Optimization of the ZGC Reference Processing Pipeline}
\author{Fredrik Hammarberg}
\authorshort{Hammarberg, F.}
\date{June 2026}
\titlepagelogo{UU_logo_sv_42.eps}
\publisher{Uppsala University}
\series{Master's Thesis in Computer Science}
\address{Department of Information Technology\\Uppsala University\\Uppsala, Sweden}

% Abstract
\abstractdummy{
	\begin{abstract}
        PLACEHOLDER
    \end{abstract}
}


\begin{document}
\frontmatter
    % Creates the front matter (title page(s), abstract, list of papers)
    % frontmatterCS = Comprehensive Summary (no half-title, good for Master's)
    % frontmatterMonograph = Monograph (includes half-title, typically for PhD)
    \frontmatterCS 
    
    % Dedication (optional)
    %\cleardoublepage
    %\vspace*{\fill}
    %\begin{center}
    %    \textit{To...}
    %\end{center}
    %\vspace*{\fill}
    
    % Acknowledgments
    \chapter*{Acknowledgments}
    \addcontentsline{toc}{chapter}{Acknowledgments}
    
    I would like to express my sincere gratitude to my supervisor for their guidance and support throughout this thesis work. It is through their help and encouragement that I was able to complete this thesis.
    
    A special thanks to the Department of Information Technology at Uppsala University for providing the resources necessary to conduct the performance evaluations and experiments described in this thesis.

	I would also like to thank my colleagues at Oracle for their insights and feedback on my work throughout the development process.
    
    Finally, I would like to thank my family and friends for their patience and encouragement throughout my studies.
    
    \vspace{1cm}
    
    Uppsala, Sweden, June 2026
    
    \vspace{0.5cm}
    
    Fredrik Hammarberg
    
    % Table of contents
    \begingroup
        \tableofcontents
    \endgroup
    
    % Lists of tables and figures
    \listoftables
    \listoffigures

\mainmatter
    % Main thesis chapters
    \chapter{Introduction}
\label{chap:introduction}
As the Java programming language moves further away from end-users and into the realm of large-scale enterprise applications, the performance of the Java Virtual Machine (JVM) becomes increasingly critical. Especially when it comes to performance on large heaps since many enterprise applications utilize large amounts of memory. The Z Garbage Collector (ZGC) is a low-latency garbage collector designed to handle large heaps with minimal pause times. However, the processing of weak references in ZGC has been identified as a potential bottleneck, particularly in applications that heavily utilize weak references. This thesis aims to analyze the performance of weak reference processing in ZGC and explore potential optimizations to improve its efficiency.


\section{Motivation}
\label{sec:motivation}



\section{Research Questions}
\label{sec:research-questions}

% TODO: List research questions

\section{Contributions}
\label{sec:contributions}

% TODO: Describe main contributions

\section{Thesis Structure}
\label{sec:thesis-structure}

% TODO: Outline thesis structure

    \chapter{Background}
\label{chap:background}

The three optimization approaches laid out in \cref{sec:implementation} rely on several key concepts and design decisions in the Java Virtual Machine's (JVM) garbage collection architecture, the Z Garbage Collector (ZGC), and the Java reference processing pipeline. This chapter provides the necessary background on these topics to understand the motivation and implementation of the optimizations. 

\section{Garbage Collection in the JVM}
\label{sec:gc-overview}

The Java Virtual Machine (JVM) provides automatic memory management through garbage collection (GC). Rather than requiring the programmer to explicitly allocate and free memory as in languages such as C or C++, the JVM tracks which objects are reachable from the application's root set comprising thread stacks, static fields, and JNI references and reclaims the memory of objects that are no longer reachable.

The JVM ships with several garbage collector implementations, each designed with different trade-offs between throughput, latency, and memory footprint:

\begin{itemize}
    \item \textbf{Serial GC}   A single-threaded, stop-the-world collector suitable for small heaps and client-style applications.
    \item \textbf{Parallel GC}   A multi-threaded throughput-oriented collector that parallelizes both young and old generation collection.
    \item \textbf{G1 GC (Garbage-First)}   A region-based collector that balances throughput and latency by incrementally collecting regions with the most garbage. G1 divides the heap into equal-sized regions and uses concurrent marking alongside evacuation pauses.
    \item \textbf{ZGC (Z Garbage Collector)}   A scalable, low-latency collector designed for heaps ranging from megabytes to terabytes. ZGC performs nearly all of its work concurrently with the application, resulting in sub-millisecond pause times regardless of heap size.
\end{itemize}

This thesis focuses on ZGC, as its concurrent design and generational architecture introduce specific challenges for reference processing that can benefit from targeted optimization.

\subsection{Generational Collection}
\label{subsec:generational}

Most modern JVM garbage collectors employ \emph{generational collection}, based on the \emph{weak generational hypothesis}: most objects die young. The heap is divided into at least two generations:

\begin{itemize}
    \item The \textbf{young generation}, where newly allocated objects reside. Young generation collections are frequent but fast, since most young objects are short-lived.
    \item The \textbf{old generation}, where objects that have survived multiple young collections are promoted. Old generation collections are less frequent but more expensive.
\end{itemize}

ZGC adopted generational collection starting with JDK~21. In generational ZGC, the young and old generations can be collected independently and concurrently. This is relevant to reference processing because references in different generations are handled during different collection cycles. Cross-generational relationships, where a reference object is in one generation and its referent is in another, require special treatment (\cref{subsec:generational-refs}).


\section{Java Reference Types}
\label{sec:reference-types}

In addition to ordinary (strong) references, Java provides a hierarchy of \emph{reference objects} in the \texttt{java.lang.ref} package that allow applications to interact with the garbage collector's reachability decisions. These reference types, listed from strongest to weakest, are:

\begin{enumerate}
    \item \textbf{SoftReference}   The referent is cleared at the discretion of the GC, typically when memory is low. Commonly used for memory-sensitive caches.
    \item \textbf{WeakReference}   The referent is cleared as soon as it becomes weakly reachable (i.e., no strong or soft references exist to it). Widely used in caching frameworks, canonical mappings, and listener registrations.
    \item \textbf{FinalReference}   An internal JVM type (not part of the public API) used to implement the \texttt{finalize()} mechanism. Objects with finalizers are wrapped in a \texttt{FinalReference} so the GC can schedule finalization before reclaiming the object.
    \item \textbf{PhantomReference}   The referent is never accessible through the reference; \texttt{get()} always returns \texttt{null}. Used for post-mortem cleanup actions, often as a replacement for finalizers.
\end{enumerate}

All four types extend the abstract class \texttt{Reference<T>} (in the \texttt{java.lang.ref} package), which defines the key fields and the lifecycle managed by the garbage collector.

\subsection{The Reference Object Lifecycle}
\label{subsec:ref-lifecycle}

Each \texttt{Reference} object maintains several fields that track its state:

\begin{itemize}
    \item \texttt{referent}   The referenced object. The GC may clear this field (set it to \texttt{null}) when the referent becomes eligible for collection according to the reference's strength.
    \item \texttt{queue}   The \texttt{ReferenceQueue} this reference is registered with. If the reference was not created with a queue, this field points to the sentinel object \texttt{ReferenceQueue.NULL\_QUEUE}.
    \item \texttt{next}   Used as a link pointer when the reference is enqueued in a \texttt{ReferenceQueue}.
    \item \texttt{discovered}   A field with dual purpose: during GC marking, it serves as a link in the garbage collector's internal \emph{discovered list}; after processing, it is reused as a link in the VM's \emph{pending list} that feeds the \texttt{ReferenceHandler} thread.
\end{itemize}

A reference object transitions through the following states during its lifetime:

\begin{enumerate}
    \item \textbf{Active}   The reference has been created and the referent is still reachable (or not yet processed). The \texttt{discovered} field may be used by the GC.
    \item \textbf{Pending}   The GC has determined the referent is eligible for clearing and has added the reference to the VM's pending list. The \texttt{ReferenceHandler} thread will process it.
    \item \textbf{Enqueued}   The \texttt{ReferenceHandler} thread has moved the reference into its associated \texttt{ReferenceQueue}, where the application can retrieve it via \texttt{poll()} or \texttt{remove()}.
    \item \textbf{Inactive}   The reference has been dequeued by the application, or was created without a queue and has been cleared. It is no longer tracked by the GC or queuing infrastructure.
\end{enumerate}

An important observation is that references created \emph{without} a \texttt{ReferenceQueue} can transition directly from Active to Inactive, bypassing the Pending and Enqueued states entirely. However, as discussed in \cref{sec:ref-processing}, the default implementation does not always take advantage of this shortcut.

\subsection{WeakReference and ReferenceQueue}
\label{subsec:weakref-queue}

A \texttt{WeakReference} is the most commonly used non-strong reference type in Java applications. It has two constructors:

\begin{figure}[htbp]
\begin{lstlisting}[language=Java]
// Without a queue - queue field set to NULL_QUEUE
public WeakReference(T referent) {
    super(referent);
}

// With a queue
public WeakReference(T referent, ReferenceQueue<? super T> q) {
    super(referent, q);
}
\end{lstlisting}
\caption{WeakReference constructors}
\label{lst:weakref-constructors}
\end{figure}

When a \texttt{WeakReference} is created with a \texttt{ReferenceQueue}, the application can be notified when the referent is collected: the GC clears the referent and eventually enqueues the reference, allowing the application to perform cleanup. When created \emph{without} a queue, the reference serves only as a conditional pointer the application can call \texttt{get()} to check whether the referent is still alive, but receives no notification upon clearing.

The \texttt{ReferenceQueue} is implemented as a synchronized singly-linked list using the \texttt{next} field of the \texttt{Reference} objects. It provides \texttt{poll()} for non-blocking retrieval and \texttt{remove()} for blocking retrieval with an optional timeout.

A key insight for this thesis is that \textbf{many real-world uses of \texttt{WeakReference} do not use a \texttt{ReferenceQueue}}. Caching frameworks (including the widely used Caffeine library~\cite{caffeine}), \texttt{WeakHashMap}, and similar data structures typically create weak references without queues, relying instead on periodic polling or lazy cleanup. This means the garbage collector expends effort preparing these references for enqueueing (adding them to the pending list) even though the \texttt{ReferenceHandler} thread will immediately skip them.


\section{The Z Garbage Collector}
\label{sec:zgc}

ZGC is a concurrent, region-based, compacting garbage collector included in the JDK since JDK~11 (experimental) and production-ready since JDK~15. Starting with JDK~21, ZGC supports generational collection, which is now the default mode. ZGC's primary design goals are:

\begin{itemize}
    \item Sub-millisecond pause times that do not increase with heap size.
    \item Support for heap sizes from 8\,MB to 16\,TB.
    \item Concurrent marking, relocation, and reference processing.
\end{itemize}

ZGC achieves these goals through several key techniques:

\paragraph{Colored Pointers.} ZGC embeds metadata in the unused bits of 64-bit object pointers (or in a dedicated area when using multi-mapped memory). These \emph{color bits} encode information about the pointer's current state relative to the GC's phase for example, whether the pointer has been remapped after relocation, or whether it points to a young or old generation object. This allows the GC to reason about pointer states without loading the object header.

\paragraph{Load Barriers.} ZGC uses a load barrier injected by the JIT compiler at every reference load. The barrier checks the color bits of the loaded pointer and, if necessary, applies a fixup (e.g., remapping after relocation or marking for concurrent marking). This is the mechanism that allows ZGC to perform most of its work concurrently with the application.

\paragraph{Concurrent Phases.} A ZGC collection cycle interleaves short stop-the-world pauses (for root scanning and synchronization) with concurrent phases (for marking, reference processing, and relocation). The concurrent phases run alongside application threads, using load barriers to maintain consistency.


\section{Reference Processing Pipeline}
\label{sec:ref-processing}

Reference processing is the mechanism by which the garbage collector handles \texttt{Reference} objects discovered during marking. The pipeline has three stages: \emph{discovery}, \emph{processing}, and \emph{enqueueing}. In ZGC, the first two stages run concurrently with the application during the marking phase, while the enqueue stage involves handing off references to the Java-level \texttt{ReferenceHandler} thread.

\subsection{Discovery}
\label{subsec:discovery}

During the concurrent marking phase, when the garbage collector encounters a \texttt{Reference} object (i.e., an instance of \texttt{SoftReference}, \texttt{WeakReference}, \texttt{FinalReference}, or \texttt{PhantomReference}), it evaluates whether the reference should be \emph{discovered} that is, added to a list for later processing.

The discovery decision depends on several factors:

\begin{enumerate}
    \item \textbf{Is the reference active?} Only active references are eligible. For non-\texttt{FinalReference} types, an inactive reference has a \texttt{null} referent. For \texttt{FinalReference}, inactivity is indicated by a non-null \texttt{next} field.
    \item \textbf{Is the referent still strongly reachable?} If the referent is already marked as strongly live, the reference does not need processing. The garbage collector checks this via its liveness information.
    \item \textbf{Soft reference policy.} For \texttt{SoftReference} objects, the collector applies a policy (based on available memory and the reference's last access time) to decide whether to keep the referent alive or clear it.
    \item \textbf{Generational considerations.} In generational ZGC, only references residing in the old generation are discovered during an old collection. References in the young generation are handled during young collections. (See \cref{subsec:generational-refs} for the implications.)
\end{enumerate}

Once a reference is selected for discovery, it is added to a \emph{discovered list}. In the standard OpenJDK implementation (used by G1, Serial, and Parallel GCs), discovered references are linked together in a per-thread singly-linked list using the \texttt{discovered} field of each \texttt{Reference} object as the link pointer. In ZGC, each GC worker thread maintains its own discovered list head, avoiding the need for atomic compare-and-swap (CAS) operations during discovery. This is in contrast to the shared reference processor (used by G1 and others), which uses CAS on the \texttt{discovered} field when multiple threads may discover references concurrently.

\subsection{Processing}
\label{subsec:processing}

After marking completes, the discovered lists are processed. For each discovered reference, the collector determines the final fate of the referent:

\begin{enumerate}
    \item If the referent has become \texttt{null} (cleared by the application concurrently), the reference is dropped from the list.
    \item If the referent is now strongly reachable (e.g., it was resurrected by another thread), the reference is dropped and the referent is kept alive.
    \item If the referent is still unreachable according to the reference's strength, the referent is \textbf{cleared} (set to \texttt{null}), and the reference is prepared for enqueueing.
\end{enumerate}

The processing order follows the strength hierarchy: soft references are processed first, then weak, then final, and finally phantom references. Processing soft references before weak references ensures that soft references are cleared only under memory pressure, while weak references are always cleared when the referent is not strongly reachable.

In ZGC, the ``clearing'' of a referent involves applying a \emph{clean barrier} that writes a null value to the referent field. Different barrier variants are used depending on the reference type: a weak-OOP variant for soft and weak references, and a phantom-OOP variant for phantom references. These barriers ensure that the pointer coloring metadata is correctly maintained.

\subsection{Enqueue}
\label{subsec:enqueue}

After a reference's referent has been cleared, the reference must be made available to the application. This happens through the \emph{pending list} mechanism:

\begin{enumerate}
    \item The GC builds an internal pending list by linking cleared references together via their \texttt{discovered} field.
    \item Under the heap lock, the GC atomically prepends its internal pending list onto the global VM pending list.
    \item The high-priority \texttt{ReferenceHandler} daemon thread is notified.
    \item The \texttt{ReferenceHandler} thread wakes up, atomically retrieves the pending list, and iterates over it. For each reference, it checks whether \texttt{queue != NULL\_QUEUE}. If the reference has a real queue, it calls \texttt{queue.enqueue(this)}; otherwise it simply advances to the next reference.
\end{enumerate}

This design reveals an inefficiency: references created without a \texttt{ReferenceQueue} are still linked into the pending list, transferred to the \texttt{ReferenceHandler} thread, and iterated over only to be skipped. For workloads where the majority of weak references lack queues, this constitutes unnecessary bookkeeping. This inefficiency is one of the key targets of the optimizations explored in this thesis.


\section{Memory Hierarchy and Pointer-Chasing}
\label{sec:memory-hierarchy}

Modern processors are equipped with deep memory hierarchies consisting of multiple levels of cache (typically L1, L2, and L3) between the CPU and main memory. Accessing data in the L1 cache takes on the order of a few nanoseconds, while a cache miss that must be serviced from main memory incurs a latency of roughly 50--100 nanoseconds, two orders of magnitude slower. The performance of any data-intensive algorithm therefore depends heavily on how well its memory access patterns interact with the caching hardware.

Two properties of memory access patterns determine cache efficiency:

\begin{itemize}
    \item \textbf{Spatial locality}: Accessing data at nearby memory addresses. When a cache line (typically 64 bytes) is loaded, all data within it becomes available at L1 speed. Sequential traversal of a contiguous array exhibits strong spatial locality, since successive elements occupy adjacent or nearby cache lines.
    \item \textbf{Temporal locality}: Accessing the same data repeatedly within a short time window. Recently accessed data remains in cache and can be re-read without penalty.
\end{itemize}

A critical distinction arises between \emph{contiguous} data structures (arrays) and \emph{pointer-linked} data structures (linked lists) when traversal involves cache misses. In a contiguous array, all element addresses are computable from the base address and index; the processor does not need to read any element in order to determine the address of the next one. Modern out-of-order processors exploit this property aggressively: the memory subsystem can issue multiple cache-miss loads simultaneously, since the addresses of future elements are known in advance. Hardware prefetchers further accelerate sequential access by detecting the stride pattern and fetching upcoming cache lines before they are explicitly requested.

In contrast, a linked list exhibits a \emph{pointer-chasing} access pattern: the address of the next node is stored inside the current node, so the processor cannot determine where the next load will go until the current load completes. Even on a wide out-of-order core capable of tracking dozens of in-flight memory operations, a linked-list traversal serializes at the rate of one cache miss per node. If each node resides on a separate cache line (as is typical for heap-allocated Java objects), the traversal stalls for the full memory latency on every step.

This difference has substantial practical consequences. When iterating over $n$ elements that all miss in cache, the array-based traversal completes in approximately one memory-latency period (since all loads can overlap), while the linked-list traversal takes $n$ times the memory latency. For the reference processing use case, where millions of \texttt{Reference} objects may be scattered across a multi-gigabyte heap, this serialization can dominate the processing time.


\section{Data Structures in Reference Processing}
\label{sec:data-structures}

The choice of data structure for storing discovered references has significant implications for performance, particularly in light of the memory hierarchy considerations described in \cref{sec:memory-hierarchy}. This section describes the structure used in the standard implementation and discusses the alternative considered in this work.

\subsection{Intrusive Linked List}
\label{subsec:linked-list}

The standard implementation stores discovered references in an \emph{intrusive singly-linked list}. Each \texttt{Reference} object's \texttt{discovered} field serves as the \texttt{next} pointer, forming a chain rooted at a per-thread list head. This design has several properties:

\begin{itemize}
    \item \textbf{Zero allocation overhead}: No auxiliary nodes are allocated; the list is embedded within the \texttt{Reference} objects themselves.
    \item \textbf{O(1) prepend}: Adding a reference to the list is a single pointer write (or a CAS in the shared processor's multi-threaded mode).
    \item \textbf{Dual-purpose field}: After processing, the same \texttt{discovered} field is reused to link references in the pending list destined for the \texttt{ReferenceHandler} thread.
\end{itemize}

However, the linked list has drawbacks relevant to modern hardware:

\begin{itemize}
    \item \textbf{Poor cache locality}: \texttt{Reference} objects are scattered across the heap. Traversing the discovered list is a pointer-chasing operation (\cref{sec:memory-hierarchy}): each load depends on the result of the previous load, serializing cache misses and preventing the processor from overlapping memory accesses.
    \item \textbf{Field overloading}: The dual use of the \texttt{discovered} field means that during processing, building the pending list requires careful coordination a reference cannot simultaneously be part of the discovered list and the pending list.
    \item \textbf{Unnecessary pending-list work}: All cleared references are threaded onto the pending list, regardless of whether they have a queue. The \texttt{ReferenceHandler} thread must then iterate the entire list to pick out the ones that actually need enqueueing.
\end{itemize}

\subsection{Dynamic Array}
\label{subsec:dynamic-array}

An alternative to the intrusive linked list is a \emph{dynamic array} that stores metadata about discovered references contiguously in memory. Rather than linking references via the \texttt{discovered} field, each entry in the array records:

\begin{itemize}
    \item A pointer to the referent field of the \texttt{Reference} object.
    \item A pointer to the discovered field.
    \item The address of the referent object.
    \item A snapshot of the referent pointer value at discovery time.
\end{itemize}

This approach has complementary trade-offs:

\begin{itemize}
    \item \textbf{Better cache locality}: Entries are stored contiguously, enabling sequential access patterns that benefit from hardware prefetching. Unlike the linked-list traversal, the processor can issue all cache-miss loads in parallel since element addresses are computable without reading previous elements (\cref{sec:memory-hierarchy}).
    \item \textbf{Direct field access}: By storing raw pointers to the \texttt{Reference} object's fields at discovery time, the processing phase can write directly to those fields without re-loading the object header or navigating the object layout avoiding one level of indirection.
    \item \textbf{Independent of the \texttt{discovered} field}: The array is a separate structure, so the \texttt{discovered} field is not overloaded. To prevent re-discovery, the field is set to a self-loop sentinel.
    \item \textbf{Allocation cost}: Unlike the intrusive list, the array requires heap allocation for its backing storage. Growth involves copying the existing entries to a new, larger buffer.
\end{itemize}

The array approach is particularly well-suited to references that will \emph{not} be enqueued, because these references never need to appear on the pending list. Their entire lifecycle from discovery through clearing can be handled within the GC's own data structures, without involving the \texttt{ReferenceHandler} thread at all.


\section{ZGC Reference Processing Details}
\label{sec:zgc-ref-details}

This section provides further detail on aspects of ZGC's reference processing that are directly relevant to the optimizations discussed in subsequent chapters.

\subsection{Per-Worker Discovered Lists}
\label{subsec:per-worker-lists}

In ZGC, each GC worker thread maintains its own discovered list head (via the \texttt{ZPerWorker} template). During concurrent marking, when a worker encounters a reference that should be discovered, it prepends it to its own list. Because each worker has an independent list head, \textbf{no atomic compare-and-swap (CAS) is required for discovery} a simple pointer write suffices.

This stands in contrast to the shared reference processor used by G1 and other collectors, where multiple GC threads may concurrently attempt to discover references. In the shared processor, an atomic CAS on the \texttt{discovered} field is used to resolve races: the CAS succeeds only if the field is still \texttt{null} (i.e., the reference has not already been discovered by another thread). While correct, this introduces contention and latency, especially under high discovery rates.

\subsection{Generational Reference Handling}
\label{subsec:generational-refs}

Generational ZGC introduces a subtlety in reference processing that has implications for correctness. When a \texttt{Reference} object resides in the old generation but its referent is in the young generation, the referent is treated as \textbf{strongly reachable} during old-generation collection. This conservative treatment ensures that the referent is not prematurely cleared it will be properly evaluated when the young generation is collected.

Conversely, only references in the old generation are discovered during an old-generation collection cycle. References in the young generation are skipped entirely and will be handled during a young-generation collection.

This generational boundary has practical consequences for the array-based optimization: when processing array entries, the GC must verify that the referent is in the old generation before clearing it. If the referent has been promoted to the young generation, it must be kept alive, and the entry is dropped without clearing. If young-generation marking is active, the referent may additionally need to be explicitly marked through the young-generation marking barrier to ensure it is not lost.

\subsection{Concurrent Processing and Barriers}
\label{subsec:barriers}

Because ZGC performs reference processing concurrently with the application, the collector must use barriers when reading or writing reference fields. The load barrier ensures that any pointer loaded by the GC is in the correct ``color'' (i.e., its metadata bits are consistent with the current GC phase). The clean barrier, used during referent clearing, writes a properly colored null value.

These barriers are essential for correctness but add overhead to each reference processing operation. The dynamic array approach stores a snapshot of the referent pointer value at discovery time, which can reduce the number of barrier operations needed during the processing phase, as the snapshot can be compared directly rather than re-applying a load barrier.

    \chapter{Method}
\label{chap:method}

In order to answer the research questions posed in \cref{sec:research-questions} three orthogonal optimisation approaches were implemented separately for the processing of weak references without a \texttt{ReferenceQueue} in the Z garbage collector (ZGC). These separate implementations were then combined in all possible combinations and had their performance evaluated using microbenchmarks. These three approaches were as presented in \cref{chap:introduction} the following:
\begin{itemize}
    \item Skipping the enqueue path for references without a \texttt{ReferenceQueue}
    \item Replacing the intrusive linked list with a dynamic array for the discovered list for references without a \texttt{ReferenceQueue}
    \item Sequential Code Optimisations
\end{itemize}
The rationale and implementation details of these approaches are described in the following sections. The evaluation of the approaches is described in \cref{chap:results}.

\section{Skipping Enqueue Path}
\label{sec:skipping-enqueue-path}
As described in \cref{subsec:processing} ZGC enqueues discovered references to the internal pending list regardless of whether they have a \texttt{ReferenceQueue} or not. This is inefficient for weak references without a \texttt{ReferenceQueue} as discussed in \cref{subsec:enqueue}. Because of this, the first optimisation approach implemented was to skip the enqueue path for references without a \texttt{ReferenceQueue}. The intended effect of this optimisation is to both reduce the amount of work done by the GC threads during reference processing and to reduce the amount of work done by the \texttt{ReferenceHandler} thread when iterating over the pending list.

\subsection{Implementation of Null-queue Fast Path}
In order to skip enqueueing references without a \texttt{ReferenceQueue} the \texttt{ReferenceProcessor} class was modified to check if the queue field of the reference being discovered is equal to the sentinel value \texttt{ReferenceQueue.NULL\_QUEUE} before adding it to the discovered list. If the reference does not have a \texttt{ReferenceQueue} it is added to a separate discovered list for references without a \texttt{ReferenceQueue} instead of being added to the regular discovered list. This separate discovered list is then processed separately by the GC threads and none of its references are enqueued to the pending list for the \texttt{ReferenceHandler} thread. This optimisation is controlled by a JVM flag and can be enabled or disabled at runtime.

\section{Replacing Intrusive Linked List with Dynamic Array}
\label{sec:replacing-intrusive-linked-list}
%Rationale and hoped-for effect of this optimisation approach

\subsection{Implementation of Dynamic Array for Queue-Less WeakReferences}

\section{Sequential Code Optimisations}
\label{sec:sequential-code-optimisations}
%Rationale and hoped-for effect of this optimisation approach

\subsection{Implementation of Sequential Code Optimisations}

\section{Performance Evaluation}
    \chapter{Results}
\label{chap:results}

This chapter presents the empirical evaluation of the implemented ZGC weak-reference optimisations and the exploratory weak-field mechanism described in \cref{chap:method}.

\section{Compared Configurations}
\label{sec:results-configurations}

For \texttt{WeakReference}-based processing, the comparison covers the baseline (\texttt{none}) and all seven non-empty combinations of the three optimisation flags (clear path, separate lists, dynamic array). In addition, a dedicated \texttt{weak\_fields} configuration is included, where weak-field benchmark classes are executed in the same harness and analysed with the same tooling.

This configuration set allows two complementary comparisons:
\begin{itemize}
	\item \textbf{Intra-model comparison} among \texttt{WeakReference} variants to isolate the impact of each optimisation and their interactions.
	\item \textbf{Inter-model comparison} between optimised \texttt{WeakReference} processing and weak-field representation to evaluate structural overhead trade-offs.
\end{itemize}

\section{Measured Metrics}
\label{sec:results-metrics}

The primary performance metrics are the GC timing statistics extracted from ZGC logs, including major collection time, old-generation collection time, concurrent mark-follow subphases, and concurrent reference-processing subphases. Memory behaviour is characterised using continuous samples of RSS, GC committed memory, and Java heap usage.

Following the methodology in \cref{subsec:measurement-methodology}, single-object benchmark comparisons use the per-run maximum values as primary statistics, while multi-object comparisons use per-run averages. This choice reflects the fact that single-object runs are designed to create concentrated peak load in one collection event, whereas multi-object runs distribute work over repeated phases.

\section{Presentation of Results}
\label{sec:results-presentation}

\Cref{fig:metric-boxplots-single,fig:metric-boxplots-multi} show box plots of the GC timing metrics for the single-object and multi-object benchmarks. They summarise the distribution of per-run maxima for the single-object benchmark and per-run averages for the multi-object benchmark, allowing comparison of both central tendency and variability. \Cref{fig:continuous-memory-single,fig:continuous-memory-multi} show representative continuous-memory traces for the same benchmark families, illustrating how memory usage evolves over time and across GC phases.

\begin{figure}
	\centering
	\includegraphics[width=\textwidth]{metric_boxplots_single_1.pdf}
	\caption[GC metric box plots (single-object)]{Box plots of GC metrics across evaluated configurations for the single-object benchmark.}
	\label{fig:metric-boxplots-single}
\end{figure}

\begin{figure}
	\centering
	\includegraphics[width=\textwidth]{metric_boxplots_double_2.pdf}
	\caption[GC metric box plots (multi-object)]{Box plots of GC metrics across evaluated configurations for the multi-object benchmark.}
	\label{fig:metric-boxplots-multi}
\end{figure}

\begin{figure}
	\centering
	\includegraphics[width=\textwidth]{continuous_memory_single_1.pdf}
	\caption[Continuous memory trace (single-object)]{Continuous memory trace for the single-object benchmark.}
	\label{fig:continuous-memory-single}
\end{figure}

\begin{figure}
	\centering
	\includegraphics[width=\textwidth]{continuous_memory_double_2.pdf}
	\caption[Continuous memory trace (multi-object)]{Continuous memory trace for the multi-object benchmark.}
	\label{fig:continuous-memory-multi}
\end{figure}

    \chapter{Discussion}
\label{chap:discussion}

\section{Interpretation of Results}
\label{sec:discussion-interpretation}

The combination of the optimised clear path and dynamic array achieves a far greater reduction in non-strong processing time than either optimisation individually. The optimised clear path eliminates the enqueue step for queue-less weak references, routing them through a lighter processing path; once that lighter path is active, the dynamic array makes traversal of the resulting list faster through improved cache locality~\cite{drepper07-memory}. Each optimisation therefore amplifies the benefit of the other, producing a super-additive interaction. The near-equivalence of \texttt{clear\_path\_dyn} and \texttt{all} follows from this: when the clear path and dynamic array are already active, adding the separate-list mechanism provides little further gain because most of the processing work has already been removed. The \texttt{sep\_only} variant confirms this from the other direction: routing references to a separate list without changing either the traversal structure or the per-reference work yields no meaningful speedup.

The \texttt{weak\_fields} variant achieves a lower non-strong processing time than the baseline but falls short of \texttt{clear\_path\_dyn} because it bypasses the \texttt{WeakReference} pipeline entirely. The residual non-strong processing time for \texttt{weak\_fields} reflects only reference types that remain in the pipeline, such as phantom and final references, rather than any optimisation of weak-reference traversal.

Despite the substantial reductions in non-strong processing time, the \texttt{WeakReference} optimisations have a limited measurable effect on total collection time. Reference processing accounts for only a fraction of the total GC cycle; the remainder is dominated by concurrent marking, relocation, and other phases that the optimisations do not touch. The overlapping distributions for major collection and old-generation time across \texttt{WeakReference} variants reflect this: once reference processing is no longer the dominant cost, further gains in that subphase translate to at most a small fraction of total cycle time.

The bimodal shape of many violin plots for overall collection metrics suggests that the benchmark encounters two distinct operating regimes. A plausible explanation is that some GC cycles carry a heavier concurrent mark-follow workload than others, perhaps due to object lifetimes or allocation patterns interacting with the timing of the explicit \texttt{System.gc()} call. This is consistent with similar bimodality observed in the mark-follow metric. The higher variability of the single-object benchmark relative to the multi-object benchmark follows from its design: concentrating all 10 million references into a single GC cycle amplifies sensitivity to transient system effects such as cache state, NUMA placement, and thread scheduling. The multi-object benchmark spreads the reference load across five cycles, averaging out per-cycle variation.

The large advantage of \texttt{weak\_fields} on overall collection metrics stems from eliminating \texttt{WeakReference} wrapper objects entirely. Each \texttt{WeakReference} object occupies heap space for its header, referent, queue, and discovered-list pointer fields, and with 10 million such objects in the single-object benchmark this amounts to approximately 460\,MB of additional live heap. Eliminating these objects reduces GC pressure across all phases: fewer objects to mark, fewer pages to scan during relocation, and less metadata to maintain. The optimised clear path and dynamic array reduce per-reference cost within the processing pipeline, but they cannot offset the broader cost of maintaining 10 million reference objects throughout the GC cycle.

The elevated auxiliary memory of \texttt{clear\_path\_dyn} (647\,MB median in the single-object benchmark) compared to \texttt{all} (455\,MB) reflects an interaction between the optimisations. Without the separate-list mechanism, all queue-less references share a single dynamic array that retains its full allocated capacity throughout the GC cycle. With the separate-list mechanism active, references are split across two independently managed lists, bounding the peak size of each more tightly. In all cases, the additional GC-committed memory is negligible relative to the 100\,GB fixed heap. The identical heap footprints of all \texttt{WeakReference} variants confirm that the optimisations do not affect the number of wrapper objects allocated; only \texttt{weak\_fields}, which eliminates the wrapper objects themselves, reduces Java heap usage.

\section{Future Work}
\label{sec:future-work}

Can the optimisations developed for ZGC be applied to other garbage collectors in the JVM, such as the G1 garbage collector, and if so, what performance improvements can be observed?
    \chapter{Conclusion}
\label{chap:conclusion}

This thesis investigated whether ZGC can process queue-less \texttt{WeakReference} objects more efficiently, and whether weak semantics expressed directly in object fields can reduce overhead further. To answer these questions, three orthogonal \texttt{WeakReference}-pipeline optimisations were implemented in an OpenJDK research fork: a queue-less fast path that skips enqueue handling, a dynamic-array representation for discovered queue-less weak references, and a specialised clear path that removes unnecessary per-reference work. In addition, an exploratory weak-field mechanism was implemented through a Java field annotation and corresponding VM support.

The evaluation shows a clear answer to \cref{rq:weakref-opt}. For the reference-processing subphase most directly targeted by the implementation, the strongest \texttt{WeakReference} variants are those that combine the specialised clear path with the dynamic array. In the single-object benchmark, \texttt{clear\_path\_dyn} and \texttt{all} reduce median non-strong processing time from 996.9\,ms to 187.9\,ms and 184.4\,ms respectively, corresponding to reductions of about 81\%. In the multi-object benchmark, the same variants reduce the median from 44.1\,ms to 18.8\,ms and 18.7\,ms, or about 57\%. The skip-enqueue separation on its own contributes little, which indicates that the main savings come from improved traversal locality and reduced per-reference work rather than from the list split alone.

For broader GC cycle metrics, the answer is more restrained. The best \texttt{WeakReference} variants reduce median major-collection time by approximately 8\% in the single-object benchmark, but those improvements are small relative to the overall spread of the measurements. In other words, optimising queue-less weak-reference processing clearly helps the targeted subphase, but it does not by itself transform the full cost profile of a GC cycle.

The answer to \cref{rq:weak-fields} is more mixed but also more structurally interesting. The weak-field mechanism does not achieve the lowest non-strong processing times; the best specialised \texttt{WeakReference} variants still perform better on that metric. However, weak fields substantially reduce Java heap usage and also improve total collection metrics, especially in the single-object benchmark. This indicates that a large share of weak-reference overhead comes from representation cost, namely the allocation and retention of millions of wrapper objects, rather than only from the mechanics of the non-strong processing phase.

Taken together, the results support two conclusions. First, ZGC's handling of queue-less \texttt{WeakReference} objects can be made materially faster with targeted, semantics-preserving implementation changes. Second, representation choices matter at least as much as pipeline mechanics: a field-based weak-semantics model can reduce broader GC costs by shrinking the heap structure that the collector has to traverse, even if its current prototype does not yet outperform the best specialised \texttt{WeakReference} path on every subphase. The thesis therefore contributes both a concrete optimisation study for ZGC's existing weak-reference pipeline and an exploratory argument that future JVM designs may benefit from supporting weak semantics more directly in object representation.

Placed in the broader context surveyed in \cref{chap:related-work}, these results also carry a wider implication. Across contemporary languages, callback-less weak semantics is the norm for the pure liveness-probing use case: Go's \texttt{weak.Pointer}~\cite{go-weak-pkg}, C++'s \texttt{std::weak\_ptr}~\cite{cppreference-weak-ptr}, .NET's \texttt{WeakReference<T>}~\cite{dotnet-weakreference}, and Lua's weak table entries~\cite{lua54-manual} all belong to this family, providing a probe mechanism without any notification callback. Java's \texttt{WeakReference} API already offers a callback-less path through its one-argument constructor, but the JVM has historically processed queue-less references identically to queue-backed ones, incurring the full cost of the notification machinery even when no notification was wanted. The queue-less optimisations developed here close that gap: ZGC can now handle references that were always conceptually equivalent to Go's \texttt{weak.Pointer} with correspondingly lower runtime overhead. The weak-field mechanism takes the convergence further, eliminating the wrapper object entirely and moving the representation closer to Lua's embedded weak entries or .NET's \texttt{ConditionalWeakTable}. Together, the two lines of work suggest a trajectory along which Java and ZGC can move towards the leaner weak-semantics designs found elsewhere in the language ecosystem, without abandoning the richer notification capabilities that distinguish Java's reference API for workloads that genuinely need them.



\backmatter
    % References
    % Using ieeetr style (modern, suitable for CS)
    % Alternative: alpha, abbrv, or acm
    \bibliographystyle{ieeetr}
    \bibliography{References}

\end{document}